\subsection*{Opgave 1}
Udregn følgende. 
\begin{align*}
	&1) \ 4(2+7)   &&2) \   \frac{6}{3}\cdot 7+3\frac{10}{2}  \\
	&3) \ 2!^3   &&4) \  (2+4)^2   \\
	&5) \  \frac{12}{4} +9 &&6) \  -5^2+9-\frac{14}{7}\cdot(-2)   \\
	&7) \ (-5)^3+\frac{24}{2+2}    &&8) \  (1+3!)^2   \\
	&9) \  (2+3)^3  &&10) \  \sqrt{3^2+4^2}   \\
	&11) \ \sqrt{(-6)^2+(-8)^2}   &&12) \  \frac{(1+1+1)!+6^2 + (2^2-6\cdot(-2))}{\sqrt{4^2+20}}   \\
\end{align*}

\subsection*{Opgave 2}

\begin{enumerate}[label=\roman*)]
	\item Udregn $(-2)(-3)$
	\item Udregn $(-5)2$.
	\item Udregn $(-3)(-0.5)(-6)(-1)$
	\item Udregn $(3)(-3)(-3)(-1)$
	\item For hvilke heltal $n>0$ gælder det, at $(-2)^n$ bliver et positivt tal?
	\item For hvilke heltal $n>0$ gælder det, at $(-2)^n$ bliver et negativt tal?
\end{enumerate}



\subsection*{Opgave 3}
Reducer følgende udtryk så meget som muligt.
\begin{align*}
	&1) \  (a+a)b  &&2) \  \sqrt{a^2}   \\
	&3) \  (a-b)a-a^2+ab+c  &&4) \ (\sqrt[7]{a+b})^7    \\
\end{align*}


\subsection*{Opgave 4}
Løs følgende ligninger.	
\begin{align*}
	&1) \  2x = 4  &&2) \ (-5+2)x+3!x = 2x   \\
	&3) \  5x+2x = 21  &&4) \ \sqrt{x+2} = 2   \\
	&5) \  (x-4)! = 24  &&6) \  2^x = 8  \\
	&7) \  3+4\cdot 2x = 11x  &&8) \ x^2 = 2x   \\
	&9) \  \frac{x^2}{3} = x  &&10) \  \frac{x+x}{2} = 1  \\	
\end{align*}


\subsection*{Opgave 5}
Udregn følgende brøker (forkort så meget som muligt). Brug Maple til at tjekke jeres svar
\begin{align*}
	&1)\  \frac{6}{7} + \frac{3}{1}            &&2)\ \frac{7}{22} + \frac{9}{10}\\
 	&3)\ \frac{4}{5} + \frac{3}{2}      &&4)\ \frac{1}{2} - \frac{3}{4}\\
 	&5)\  2\frac{2}{3} + \frac{7+4}{2}              &&6)\ \frac{4}{\frac{5}{7}}\\
 	&7)\  \frac{\frac{2}{3}}{6} - 2           &&8)\ \frac{-7+\frac{2}{6}}{8} + \frac{9}{5}\\
	&9)\ \frac{\frac{4}{3}}{\frac{2}{3}} &&10)\ \frac{\frac{10}{3}-\frac{2}{4}}{\frac{4}{3}+\frac{5}{3}} \\
	&11)\  \frac{\frac{1}{2}-\frac{7}{10}}{\frac{2}{5}+\frac{11}{3}} -  \frac{\frac{22}{3}+ \frac{-23}{6}}{\frac{1}{2}}        &&12)\ \sqrt{\frac{16}{25}} + \left(\frac{3}{2+4}\right)^2\\
	&13)\ \sqrt{\frac{\frac{100}{36}}{\frac{25}{49}}} && 14) \   \left(\frac{\sqrt{\frac{5}{7}}}{\sqrt{\frac{2}{9}}}\right)^2 + \sqrt[3]{\frac{\left(\frac{2+5}{3+11}\right)^3}{\left( \frac{5}{7-6}\right)^3}}
\end{align*}

\subsection*{Opgave 6}
Forkort følgende brøker så meget som muligt. 
\begin{align*}
	&1) \ \frac{ab}{a}   &2) \  \frac{a+b}{b} + \frac{a-c}{a}    \\
	&3) \ \frac{(a+b)b-ab}{b}   &4) \  \left(\frac{a}{b}\right)^2 - \frac{a^3}{b^3}   \\
\end{align*}

\subsection*{Opgave 7}
Løs følgende ligninger.
\begin{align*}
	&1) \ \frac{x}{9} = \frac{4}{x}   &&2) \ \frac{x}{4+\frac{2}{5}} = 2     \\
	&3) \ \sqrt{\frac{x}{4}} = \frac{\frac{2}{5}+\frac{16}{10}}{\frac{1}{2}}    &&4) 
	\ \frac{\frac{1}{4}}{\frac{1}{8}} + \frac{x+\frac{4}{2}x}{2} = \frac{7}{\frac{1}{2}}    \\
\end{align*}

\subsection*{Opgave 8}
Løs følgende ligninger
\begin{align*}
	&1) \ \ \frac{-7x-12}{6} + \frac{2x + 12}{6} - \frac{1}{3} = \frac{-8x + 1}{14} + \frac{x-9}{12} \\
	&2) \ \ \frac{5x}{6} - \frac{-2x-2}{3} = \frac{x + 5}{12} + \frac{4x - 3}{3} \\
	&3) \ \ \frac{-7x + 2}{9} + \frac{5x-8}{4} = \frac{-5x + 8}{20} + \frac{-x - 9}{15} \\
	&4) \ \ \frac{5x}{8} + \frac{-9x+7}{28} = \frac{8x + 6}{28} + \frac{2x + 2}{30} + \frac{3}{4} \\
	&5) \ \ \frac{-8x}{18} + \frac{3x + 1}{28} = \frac{-3x + 12}{28} - \frac{3x-10}{14} - \frac{5}{4} \\
\end{align*}
